% Open Questions — LaTeX section for arXiv:1703.05169 replication
\section*{Open Questions (5)}

\begin{enumerate}

\item \textbf{Is the Heisenberg-scaling claim ensemble-robust, or just a
    single-run assertion?}
  \emph{Basis.}  Current work reports one representative-run ratio
  ($1.12\times$ the $1/N_{\text{tot}}$ bound at seed~38, $\sum M=5981$).  The
  200-trial scaling ensemble was used only to show the plateau/collapse
  shape; no log--log post-lock slope was fit and no ensemble-mean
  $\sigma\cdot N_{\text{tot}}$ constancy was reported.  With a
  heavy-tailed run-to-run distribution (Experiment~C: 4\%
  below-headline at 50 steps), a single-run bound-ratio can overstate
  saturation.
  \emph{Next steps.}  Reuse \texttt{experimentB\_scaling.json}: (i) per-trial
  detect post-lock step $k^\star$ (first $M>1$); (ii) fit $\log\text{err}$
  vs $\log\sum M$ over $k\in[k^\star,50]$, report median slope with
  25/75-percentile band; (iii) compute $\sigma_{\text{step}}\cdot
  N_{\text{tot}}(\text{step})$ per trial and plot ensemble median with
  IQR --- a Heisenberg-saturated protocol keeps it constant.  No new sims.

\item \textbf{Does our Qiskit-shot replication reproduce the paper's own
    1000-run classical-sim dashed curve on a shot-for-shot basis?}
  \emph{Basis.}  The paper's Fig.~2a plots one representative chip run and a
  separate 1000-run classical average.  We match/beat the single-run
  number ($1.87\times 10^{-4}$ vs $2.4\times 10^{-4}$~rad) but have a
  posterior $\sigma$ that is $3\times$ larger --- suggesting the
  20000-particle Gaussian-refit and the paper's stated 1000-particle SMC
  handle the tails differently.  A shot-for-shot classical-vs-classical
  overlay would isolate chip-noise from algorithm-noise.
  \emph{Next steps.}  Rerun Experiment~A at $n_p=1000$ across 1000 seeds,
  average per-step median-err, and overlay on the paper's Fig.~2a dashed
  curve (extract with WebPlotDigitizer).  Overlap within 20\%
  median-err at every step confirms algorithm equivalence and localizes
  the paper's chip-vs-sim gap to hardware noise.  Free endpoint,
  $\approx 15$~min wallclock.

\item \textbf{Would the replication survive a realistic silicon-photonic
    noise model (SFWM heralding inefficiency, SNSPD dark counts,
    thermo-optic phase drift, gate infidelity)?}
  \emph{Basis.}  C4 was deferred; the current code uses the ideal
  Qiskit statevector (perfect gates, no depolarizing $T_2$, no heralding
  loss).  The paper's headline noise claims are tolerance up to
  $\sigma_\phi\!\approx\! 0.3$~rad gate phase noise and polynomial
  (not exponential) degradation under depolarizing $T_2$.  Without a
  noise sweep, the current statevector-perfect result is a best-case
  ceiling.
  \emph{Next steps.}  Extend \texttt{rfpe\_sim.py} with three toggles:
  (i) additive Gaussian jitter on $\Theta$ at
  $\sigma_\phi\in\{0,0.05,0.1,0.2,0.3,0.5\}$~rad; (ii) per-step
  depolarizing channel with $p=1-\exp(-M/T_2)$,
  $T_2\in\{\infty,100,30,10,3\}$; (iii) Bernoulli heralding with
  $\eta\in\{1.0,0.8,0.5,0.2\}$.  Repeat the 200-trial ensemble on each
  grid point.  Test whether the median final-err curve matches
  Fig.~3 of the paper.  Free endpoint, $\sim 1$--$2$~h wallclock.

\item \textbf{Does the RFPE + Ferrie particle-guess heuristic generalize
    to a multi-parameter Bayesian QPE (jointly estimating two commuting
    phases $\Phi_1,\Phi_2$ on a controlled-$U(4)$) with the same
    per-parameter Heisenberg scaling?}
  \emph{Basis.}  The paper's algorithm is scalar (one $\Phi$, one Gaussian
  prior).  Rudinger et al.~2017 and Gebhart et al.~2023 propose
  multi-parameter Bayesian QPE using joint SMC on a product-of-Gaussians
  prior; whether the Ferrie heuristic $M=\lceil 1.25/\sigma\rceil$,
  $\Theta\sim P(\phi)$ generalizes cleanly to a vector prior with the
  same convergence rate is an open experimental question.
  \emph{Next steps.}  Modify \texttt{rfpe\_sim.py}: replace $\Phi$ with a
  2-vector $(\Phi_1,\Phi_2)$ and $U$ with
  $\operatorname{diag}(1,e^{i\Phi_1},e^{i\Phi_2},e^{i(\Phi_1+\Phi_2)})$ on a
  3-qubit circuit; use a diagonal-covariance Gaussian prior; pick $(M_i,
  \Theta_i)$ independently per parameter; refit a joint Gaussian per
  step.  Measure per-parameter final err at 50, 100, 200 steps and check
  Heisenberg scaling holds on each axis.  Free endpoint.

\item \textbf{Does swapping the Gaussian SMC prior for a matrix-product-state
    (tensor-network) prior fix the low-$M$ plateau failure mode
    ($\approx 73\%$ of seeds in Experiment~C got stuck at 50 steps)?}
  \emph{Basis.}  Experiment~C revealed a bimodal outcome: most runs either
  escape the $M=1$ plateau and collapse super-exponentially, or stay
  stuck near the initial prior.  This is a known SMC pathology when the
  Gaussian is a poor approximation to a truly bimodal / long-tailed
  posterior.  Tensor-network / MPS priors (Stoudenmire \& Schwab 2016;
  Cheng et al.~2019) natively represent multi-modal distributions and
  may cure the plateau without needing $n_p=20000$.
  \emph{Next steps.}  Reimplement the prior as a 1-D MPS with bond
  dimension $\chi\in\{2,4,8,16\}$ on a 1024-bin phase discretization;
  keep the Ferrie heuristic for $(M,\Theta)$; replace resample--reweight--
  refit with an MPS Bayesian update (Hadamard-product of the
  likelihood-as-vector, then compress).  Rerun Experiment~C on 100
  seeds; compare (i) median final err, (ii) fraction below $10^{-2}$~rad,
  (iii) stuck fraction.  Predicted: at $\chi=8$ the stuck fraction drops
  from $\sim 73\%$ toward $\sim 10\%$.  Free endpoint; use \texttt{quimb}
  or \texttt{tenpy}.

\end{enumerate}
